% ============================================================== PROTOCOL ==
\section{Information boundaries and research protocol}\label{sec:protocol}

\subsection{What each task may read}
Tasks 1, 3 and 4 are offline: the rules allow any released observation of the evaluation split, at any time,
including times after the target cell. Task 2 is online: a forecast for origin $T$ may use released data stamped at
or before $T$. We adopted a stricter contract for Task 2 and kept it in every selected file
(\cref{tab:infosets,fig:t2window}). A forecast reads \emph{exactly} the twelve released slots $T{-}60,\dots,T{-}5$
of its own window, static network data, and models or analogue libraries fitted on the training split. It does not
read slot $T$ (which the current rules allow), other days or windows of the evaluation month, or the offline state
reconstruction. The organizer's own persistence baseline was rewritten to the same twelve slots during the
competition.

\begin{figure}[h]
\centering
\begin{tikzpicture}[x=6.3mm,y=7mm,font=\sffamily\scriptsize]
  % history slots
  \foreach \i in {0,...,11}{
    \fill[accent!78] (\i,0) rectangle ++(0.86,0.86);
  }
  \fill[gridc] (12,0) rectangle ++(0.86,0.86);
  \foreach \i in {13,...,18}{
    \fill[orange!80] (\i,0) rectangle ++(0.86,0.86);
  }
  \foreach \i in {19,...,22}{
    \fill[gridc!70] (\i,0) rectangle ++(0.86,0.86);
  }
  \node[text=white] at (0.43,0.43) {\tiny $-60$};
  \node[text=white] at (11.43,0.43) {\tiny $-5$};
  \node at (12.43,0.43) {\tiny $T$};
  \node[text=white] at (13.43,0.43) {\tiny $+5$};
  \node[text=white] at (18.43,0.43) {\tiny $+30$};
  \draw[decorate,decoration={brace,amplitude=4pt}] (0,1.05) -- (11.86,1.05) node[midway,above=4pt,align=center]
     {\textbf{Read:} 12 released slots, $T{-}60 \dots T{-}5$\\speed, flow, occupancy, ramps, eligibility};
  \draw[decorate,decoration={brace,amplitude=4pt}] (13,1.05) -- (18.86,1.05) node[midway,above=4pt,align=center]
     {\textbf{Predict:} queue field\\$L\times 6$ binary cells};
  \draw[muted,line width=0.4pt] (12.43,-0.12) -- (12.43,-0.42);
  \node[below,align=center,inner sep=1pt] at (12.43,-0.42) {slot $T$: released,\\\emph{not read}};
  \draw[decorate,decoration={brace,amplitude=4pt,mirror}] (19,-0.18) -- (22.86,-0.18) node[midway,below=4pt,align=center]
     {blanked buffer;\\later data never read};
  \node[below,align=center,inner sep=1pt] at (5.43,-0.3)
     {\textbf{+} static network, TRAIN-fitted models,\\TRAIN analogue library};
\end{tikzpicture}
\caption{The Task 2 information contract used in every selected file. Each window reads only its own twelve legal
history slots. Other evaluation days, other windows, slot $T$, ramp values after $T$ and the offline Task 1
reconstruction are excluded.}\label{fig:t2window}
\end{figure}

\begin{table}[h]
\centering\small
\caption{Information sets. ``Allowed'' is the competition rule; ``used'' is what the selected files read.}\label{tab:infosets}
\begin{tabularx}{\linewidth}{@{}l X X@{}}
\toprule
Task & Allowed by the rules & Used by TF3 and ES4 \\
\midrule
1 State & Any released observation of the split, any time; train in full & Masked layer of the same split (target channels absent by construction); TRAIN-fitted trees, profiles, factor and covariance models \\
2 Queue & Released data stamped $\le T$, including earlier days of the split & Exactly $T{-}60\dots T{-}5$ of the window; static network; TRAIN-fitted classifier and TRAIN analogue library \\
3 Physics & No rows; read from Task 1 & --- \\
4 ODME & Split's released counts, prior and incidence & Same; each split solved on its own counts and prior, never TRAIN counts \\
\bottomrule
\end{tabularx}
\end{table}

One early file (\code{P-slotT}, 30 September) set ongoing queues from the state at slot $T$. It was a diagnostic
to measure what slot $T$ adds; it was never a candidate and is not part of either selected file. \Cref{sec:ongoing}
shows that a legal forward-fill of the twelfth slot scores higher than that diagnostic on the public board.

\subsection{A hierarchy of proxies}
Hidden labels exist only on Kaggle, so every decision used one of four local instruments, each with a known weakness.
\begin{enumerate}[label=(\roman*)]
  \item \textbf{Held-out training months.} Models are fitted on June--November and scored on December, January and
        February with the official masks and the official scorers for Tasks 1 and 4. This is the cleanest instrument,
        but training months share one generator draw: they are not new incident schedules.
  \item \textbf{The 80 published training queue windows} (40 onset, 40 ongoing), all dated 1--10 June 2030, scored
        with latent labels reconstructed from the unmasked layer. They are few and repeatedly inspected.
  \item \textbf{Visible-cell pseudo-events in March and April.} Extra masks or pseudo-windows are drawn on cells the
        release shows, with odd and even dates split into development and confirmation halves. The targets are noisy
        observations, not organizer latent labels.
  \item \textbf{A local conservation proxy} that replaces the withheld boundary flux by the observed accumulation
        change. We verified that our implementation equals the official Task 3 scorer given the same flux file, but
        the proxy sits near the exact-answer ceiling and is not the leaderboard physics score.
\end{enumerate}

\subsection{Pre-registered gates}
From 3 October onward each candidate change was registered before its result was seen, with a fixed development
rule, a single confirmation and no rescue. The state ensemble (TS3, \cref{sec:ensemble}) illustrates the template:
\begin{itemize}
  \item \emph{Development} on December and January with models fitted on June--November: the joint proxy
        $0.35\,\Sstate + 0.15\,S_{\mathrm{phys}}^{\mathrm{proxy}}$ must rise by at least 0.0001, $\Sstate$ must not fall,
        at least seven of ten panels must improve and no panel may lose more than 0.001.
  \item \emph{One confirmation} on February with the same rule and at least eight improving panels. A failed arm is
        closed; no member, weight or panel subset is re-selected afterwards.
  \item \emph{Final refit} of the frozen rule on all nine training months, with unchanged seeds and settings.
\end{itemize}
Queue changes used the analogous structure on visible March/April events (odd dates for development, even dates
for confirmation) and on the published June windows, with date-clustered bootstrap intervals.

\subsection{Release engineering}
Every candidate from RC1 (2 October) to TF3 (4 October) passed the checklist below before its single upload.
The later ES files change one column of a frozen, hash-verified parent; each was checked by recomputing its
SHA256 and by a row-level difference against that parent (for example, ES2 changes exactly 1,715,267 speeds and
nothing else). Every file was uploaded once.
\begin{enumerate}
  \item An independent standard-library parser reads all 6,980,503 rows and checks keys, value domains, horizon
        placement and that only the intended columns changed relative to the parent file.
  \item A clean regeneration from the raw release and the saved weights, in a fresh pinned environment that rebuilds
        every cache, must reproduce every queue label exactly and every continuous value to within $10^{-8}$.
  \item The archive SHA256 is recorded and an exclusive attempt marker is created before the single upload call. An
        ambiguous network outcome is never retried; scores are confirmed by a separate read-only listing.
\end{enumerate}
